% ============================================================
%  Section 9 — Conclusion
% ============================================================

\section{Conclusion}
\label{sec:conclusion}

This paper presented the first systematic literature review dedicated to
\emph{repository representation paradigms for LLM-based agents}, covering
163~primary studies published between 2020 and 2025.  We identified seven
paradigms (PAR-1 to PAR-7), mapped them to five SE objectives, and derived a
multi-dimensional trade-off framework for principled paradigm selection.

Our analysis reveals that no single paradigm dominates across all dimensions:
dense embeddings offer scalability; graph-based approaches offer structural
fidelity; agentic exploration offers freshness at the cost of latency.
Matching paradigm to deployment constraints—guided by Table~\ref{tab:tradeoffs}
—is therefore essential.

We identified five open research gaps, the most critical being the absence of
multi-paradigm evaluation benchmarks and the stale-context bias demonstrated
by HumanEvo~\cite{zheng2025humanevo}.  Addressing these gaps is paramount for
advancing the field.

All supporting artefacts—snowballing log, data-extraction forms, synthesis
spreadsheets, and analysis scripts—are archived in the online replication
package on Zenodo (DOI to be assigned upon acceptance), consistent with ACM
TOSEM replication standards.
